泰勒级数将解析函数在某一点附近表示为无穷幂级数，前提是该级数在相应邻域内收敛到这个函数。截断该级数，就得到函数的多项式近似。形式化地，函数的泰勒级数由其各阶导数在同一点处的取值构成无穷级数。给定函数 $f(X)$，其以 $a$ 为中心的泰勒级数表示为：

$
f(a) + \dfrac{f'(a)}{1!}(X - a) + \dfrac{f''(a)}{2!}(X - a)^2 + \dfrac{f'''(a)}{3!}(X - a)^3 + \cdots = \sum\limits_{d=0}^{\infty}\dfrac{f^{(d)}(a)}{d!}(X - a)^d
$

截断到次数 $D$ 得到第 $D$ 个泰勒多项式 $T_D(X)$。在泰勒级数表示 $f$ 的邻域内，误差在 $D\to\infty$ 时趋于零，但它不必在每次增加次数时都减小，也不必随到 $a$ 的距离单调变化。若实函数在连接 $a$ 与 $X$ 的区间上具有 $D+1$ 阶连续导数，泰勒定理给出
\[
|f(X)-T_D(X)|\le\frac{\sup_{u}|f^{(D+1)}(u)|}{(D+1)!}|X-a|^{D+1},
\]
其中上确界取遍该区间。在级数的收敛区域之外，提高次数一般不能保证精度改善。